Commentaires sur votre proposition (s’il y en a):


L’affirmation selon laquelle les études récentes tendraient à délaisser les méthodes non supervisées au profit d’approches supervisées mérite d’être nuancée. Si certaines techniques de clustering non supervisé semblent effectivement moins mobilisées aujourd’hui, cela ne saurait être généralisé à l’ensemble des méthodes non supervisées. Des approches telles que le BDI, qui n’est pas à proprement parler supervisé, connaissent au contraire un regain d’intérêt, comme en témoigne notamment CHR2025. Par ailleurs, la prédilection pour les SVM apparaît largement contextuelle et semble relever davantage d’une tradition francophone que d’une appropriation méthodologique plus large.


En outre, les méthodes non supervisées comme BDI ont déjà fait la preuve de leur pertinence dans des domaines exigeants tels que la poésie, avec des résultats jugés convaincants, ce qui invite à relativiser leur supposé désengagement.


Par ailleurs, certaines limites des approches supervisées, en particulier des SVM, ne sont pas véritablement interrogées ici. La question de la généralisation à d’autres auteurs reste ouverte : que se passe-t-il lorsque l’on introduit un auteur supplémentaire ? SVM est toujours mené à la décision, et c'est un problème...


L’analyse s'opère enfin à l’échelle du recueil dans son ensemble, plutôt qu’au niveau de chaque pièce prise individuellement. Une telle démarche gagnerait à être validée par des expérimentations menées sur plusieurs autres corpus d'auteurs reconnus mêlant également prose et vers, afin de montrer si une telle approche fonctionne, avant de l'appliquer à son propre corpus.


Enfin, les résultats obtenus pour Beaulieu demeurent globalement négatifs. Dans ces conditions, proposer cet auteur comme candidat potentiel apparaît pour le moins hasardeux. Une conclusion plus prudente — admettant un résultat négatif ou une absence de signal significatif — serait non seulement méthodologiquement recevable, mais sans doute plus conforme aux résultats produits, au regard du travail engagé.


********************************


Cet article long montre, à l'aide d'une analyse stylométrique basée sur un modèle SVM appris de façon supervisé, que la paternité de l'œuvre "la Chanson Spirituelle" (1545) peut être envisagée de façon plus nuancée qu'il n'était admis jusqu'à présent.
L'argumentaire, clair et bien étayé, montre l'opportunité d'appliquer les dernières avancées dans les outils de normalisation automatique de la langue aux textes en ancien ou moyen français, et montre clairement le caractère innovant des travaux.
Le reste de la chaîne de traitement semble relativement classique, même si le recours à des modèles supervisés semble encore marginal.
L'approche générale est toutefois menée de façon convaincante et les analyse quantitatives et qualitatives proposées sont intéressantes.

Ma seule inquiétude réside dans l'existence de biais potentiels dans les corpus d'entraînement, qui au-delà de leur taille réduite, peuvent induire par leur sélection même une forte capacité à influencer les résultats. La justification des éléments sélectionnés pourrait être davantage détaillée, de même que l'approche par validation croisée pourrait s'intéresser à tester la robustesse des conclusions tirées en fonction des échantillonnages réalisés dans les données d'entraînement (en considérant les œuvres séparément par exemple).

Quelques petites précisions pourraient également être apportées :
- Les transcriptions automatiques des images sources ont-elles été corrigées manuellement ?
- Serait-il possible d'ajouter dans la figure 3, en abscisse, une indication des différentes parties de l'œuvre concernée (pour faire écho à la discussion de la section 5) ?



********************************